\subsection{The two profile functionals}\label{sec:profile-data}
Write \(z=(r,q)\in[0,1]^2\). The limiting density \(\nu(z)\) has uniform row and column marginals.
The diagonal coordinates are \(d=r-q\) and \(a=r+q\), and we integrate along their fibers with
respect to \(dr\). If \(m_f(y)\) is the integral of \(\nu\) along the line with coordinate \(y\)
in the family \(f=D,A\), its final vacancy is \(u_f(y)=1-m_f(y)\). For the residual fraction \(x\)
we set
\[
 U_f(y;x)=u_f(y)+x(1-u_f(y)),\qquad U_R=U_C=x,\qquad
 W(z)=\frac{\nu(z)}{U_D(d;x)U_A(a;x)} .
\]
\paragraph{Rook correction.}
We scale \(v_x(z)=U_D(d;x)U_A(a;x)\) to give it marginals of mass one, obtaining
\(p_x(r,q)=s_x(r)v_x(r,q)s_x(q)\). Symmetry gives the same factor on both sides.
Multiplying \(v_x\) by the separable factor \(\nu/(u_Du_A)\) gives the same scaled kernel.
At \(x=1\) the scaled kernel equals one. Let \(B(r,q)=p_x(r,q)-1\) be a kernel on \(L^2[0,1]\).
Operator products are defined by integration. We set
\[
 T=(I-B^2)^{-1}-I,\qquad H=TB,\qquad
 j(r)=(T-B^2)(r,r),\qquad j_0=\int_0^1j(r)\,dr .
\]
Here \(\langle f\rangle=\int_{[0,1]^2} f(r,q)\,dr\,dq\).
Powers inside these averages are pointwise powers of kernels.
\begin{align*}
 F(x)={}&\tfrac12 j_0^2+\operatorname{tr}(T^2)-\tfrac32j_0-\tfrac32\int_0^1 j(r)^2\,dr\\
 &+\langle j(r)(B+H-T)(r,q)j(q)\rangle\\
 &-\tfrac23\langle T^3\rangle+2\langle B H^2\rangle+\tfrac23\langle H^3\rangle .
\end{align*}
In particular, \(B=0\) at \(x=1\), giving \(F(1)=0\).
\paragraph{Line prefactor.}
For \(x>0\), use a midpoint grid with an even number \(m\) of points per side, and let \(J(z)\)
denote the four lines through square \(z\). There are \(m,m,2m-1,2m-1\) lines of types
\(R,C,D,A\). Set \(\rho_l=1/m\) on rows and columns and \(\rho_l=U_l/(xm)\) on diagonals of
both families. At the intersection \(z\) of two distinct lines, set
\[
 K_{lh}=W(z)\prod_{a\in J(z)\setminus\{l,h\}}U_a,\qquad
 v_R=1,\quad v_C=-1,\quad v_D=v_A=0,\qquad \Pi=vv^{\mathsf T}/2.
\]
The other entries of \(K\) are zero. A diagonal and an antidiagonal intersect on the grid only
when their parities agree. This gives rise to the two sheets of each diagonal family in the continuum.
Writing \(\rho\) also for the diagonal matrix, we define
\begin{align*}
 L_0&=K(I+\rho K+\rho\Pi)^{-1}\rho K-3\Pi,\qquad
 \eta_l=\begin{cases}1,&l\in R,C,\\ x/U_l,&l\in D,A,\end{cases}\\
 s_l&=m^{-1}\sum_{z\in l}\nu(z)W(z),\qquad
 \theta_l=3\eta_l s_l-\sum_h K_{lh}^2\rho_h .
\end{align*}
Let \([f]_m=m^{-2}\sum_z\nu(z)f(z)\). The continuum line prefactor is the limit of
\begin{align*}
 \mathcal V_m(x)={}&4[W^2]_m
 -\tfrac32\Big[W\sum_{l,h\in J(z)}(L_0)_{lh}\Big]_m\\
 &-\tfrac32\Big[W\sum_{l\in J(z)}\theta_l\Big]_m
 +\frac{9}{4m}\sum_l\eta_l s_l^2 .
\end{align*}
The sum over \(l,h\) runs over ordered pairs and includes the diagonal. The value at zero is
obtained by continuation from \(x>0\). Appendix~\ref{app:interior} identifies the insertion
terms of \(d(x)\) using this same profile and the same parity convention.
